\chapter{Motor- und Akkumodelle}
\label{ch:motor-akku}

\section{Motor}

Die Klasse \texttt{Motor} bildet eine einfache lineare Umrechnung vom Raddrehmoment auf den Strom ab. Verwendet wird die im Projekt explizit geforderte Motorkonstante:
\begin{equation}
I = \frac{M}{k_{\mathrm{m}}}
\end{equation}
Die Motorkonstante ist in \si{\newton\metre\per\ampere} angegeben. Der Motor ist als eigene Klasse modelliert, damit die Beziehung zwischen mechanischer Last und elektrischem Strom explizit und austauschbar bleibt.

\section{Abstrakte Akku-Schnittstelle}

\texttt{BatteryBase} definiert nur das Interface. Jede konkrete Akku-Klasse muss mindestens die Aktualisierung des Ladezustands und die Berechnung der Klemmenspannung bereitstellen. Diese Abstraktion ist entscheidend fuer den Polymorphismus im Simulator: Der Code arbeitet auf dem Basistyp, nicht auf einer konkreten Unterklasse.

\section{Gemeinsame Akku-Funktionalitaet}

\texttt{BatteryPack} implementiert ein Ersatzschaltbild aus Leerlaufspannung und Innenwiderstand. Der Ladezustand wird ueber die Stromintegration aktualisiert:
\begin{equation}
\mathrm{SoC}_{k+1} = \mathrm{SoC}_k - \frac{I\,\Delta t}{C_{\mathrm{nom}}}
\end{equation}
Die Nennkapazitaet wird von \si{\ampere\hour} nach \si{\ampere\second} umgerechnet. Positiver Strom entlaedt den Akku, negativer Strom laedt ihn. Die Klemmenspannung ergibt sich als
\begin{equation}
U = U_{\mathrm{OC}}(\mathrm{SoC}) - R_{\mathrm{int}} I
\end{equation}
und wird temperaturkorrigiert, wenn eine Temperatur gesetzt wurde.

Temperatur beeinflusst zwei Groessen: Der Innenwiderstand steigt bei kaelteren Temperaturen relativ zur Referenz von \SI{25}{\celsius}, waehrend die nutzbare Kapazitaet sinkt. Beide Korrekturen sind nach unten begrenzt, damit bei extremen Temperaturen keine unrealistische Extrapolation entsteht.

\section{LiPo und NMC}

LiPo- und NMC-Akkus erben beide von \texttt{BatteryPack} und unterscheiden sich nur in der Kennlinie und im Innenwiderstand. Das Projekt verwendet tabellarische Leerlaufspannungskennlinien mit linearer Interpolation \cite{numpy_docs}. Dadurch laesst sich das Verhalten beider Akkus unter identischem Lastprofil direkt vergleichen.

\begin{table}[htbp]
\centering
\caption{Gemeinsamkeiten und Unterschiede der beiden Akkupacks}
\label{tab:lipo-nmc}
\begin{tabularx}{\textwidth}{@{}lXX@{}}
\toprule
Aspekt & LiPoBatteryPack & NMCBatteryPack \\
\midrule
Vererbung & Erbt von \texttt{BatteryPack} & Erbt von \texttt{BatteryPack} \\
Serienschaltung & 10 Zellen & 10 Zellen \\
Innenwiderstand & \SI{8}{\milli\ohm} pro Zelle, skaliert mit \texttt{n\_parallel} & \SI{7}{\milli\ohm} pro Zelle, skaliert mit \texttt{n\_parallel} \\
OCV-Kennlinie & LiPo-Stuetzstellen aus Tabelle & NMC-Stuetzstellen aus Tabelle \\
Interpolation & \texttt{numpy.interp} & \texttt{numpy.interp} \\
Verhalten unter Last & Hoehere Leerlaufspannung bei gleichem SoC & Niedrigere Leerlaufspannung bei gleichem SoC \\
\bottomrule
\end{tabularx}
\end{table}

\section{Bremswiderstand und Rekuperation}

Wenn beim Bergabfahren oder Bremsen ein negativer Motorstrom entsteht, steigt der Ladezustand des Akkus. Ueberschreitet dieser Ladevorgang 100\,\%, wird der zulaessige Strom begrenzt. Der ueberschuessige Strom wird dann ueber den Bremswiderstand in Waerme umgesetzt und in \texttt{letzte\_dissipierte\_leistung\_W} gespeichert. Im aktuellen Datensatz tritt dieser Fall nicht auf; die gespeicherte dissipierte Leistung bleibt in beiden Simulationslaeufen bei null.
